\subsection{Otto equation gate and Szilassi closure tick}
\label{sec:bt1441-bt1443-otto-szilassi-closure}

The Otto import gate now blocks equations (49), (50), (64), (65), and (66) until
their rendered equation bodies are manually transcribed.  Equation (65) is the
active target because the paper states that it uses a power of the ratio of 12
icosahedral strands to 13 half-turns.

The odd half-turn problem admits a toroidal closure candidate.  Otto's path has
\[
13\pi=12\pi+\pi,
\]
so twelve half-turns can be paired while one half-turn remains as a closure tick.
The Szilassi face action supplies exactly the needed orbit type:
\[
7=2+2+2+1.
\]
The three paired face-orbits absorb the paired part; the unique bilaterally
symmetric hexagon is the candidate carrier for the odd closure tick.  This is
dual to the BT803 seven-realization census
\[
5\text{ Csaszar}+2\text{ Szilassi}=7,\qquad E=21,\qquad 42=6\cdot7,
\]
where every realization keeps a C2 geometric symmetry.

The ordered incidence lift is
\[
12(13+1)=168=21\cdot8.
\]
It maps the 12-by-14 Otto active bins into 21 Fano flags with eight local states
per flag.  The current morphism is a verified bijection but not yet canonical;
the next proof must use the Csaszar/Szilassi C2 axes or the seven Frobenius
involutions to fix the ordering.
